\documentclass{article}
\usepackage{iclr2025_conference,times}

\input{math_commands.tex}

\usepackage{hyperref}
\usepackage{url}
\usepackage{graphicx}
\usepackage{booktabs}
\usepackage{tabularx}
\usepackage{array}
\usepackage{enumitem}
\usepackage{algorithm}
\usepackage{algorithmic}
\usepackage{xcolor}

\definecolor{blockblue}{RGB}{245,248,252}
\definecolor{linegray}{RGB}{120,120,120}

\newcolumntype{L}[1]{>{\raggedright\arraybackslash}p{#1}}

\title{AI Director Agent: Explainable Chinese Storyboard Generation\\[0.25em] with Hybrid Planning and Director-Level Image Control}

\author{Lee Zi Zuo \And Thai Tian Qi \And Carole Hew \\
College of Computer Science and Artificial Intelligence \\
Fudan University \\
Shanghai, China
}

\iclrfinalcopy

\makeatletter
\def\@maketitle{%
  \vbox{\hsize\textwidth
    \centering
    {\LARGE\sc \@title\par}
    \vskip 1.4em
    {\large\bfseries
      \begin{tabular}{c@{\hspace{4.5em}}c@{\hspace{4.5em}}c}
        Lee Zi Zuo & Thai Tian Qi & Carole Hew \\
      \end{tabular}\par}
    \vskip 0.8em
    {\normalsize
      College of Computer Science and Artificial Intelligence\\
      Fudan University\\
      Shanghai, China\par}
    \vskip 0.35in minus 0.1in
  }%
}
\makeatother

\begin{document}

\maketitle

\begin{abstract}
We present \textbf{AI Director Agent}, an end-to-end system that converts short Chinese screenplays into cinematic storyboard panels with explainable intermediate reasoning.
Unlike single-prompt text-to-image tools, our pipeline decomposes narrative text into semantically parsed shots, applies cinematography rules, refines diffusion prompts with a large language model (LLM), generates images through a pluggable backend, and composes a horizontal storyboard for presentation.
A key contribution is a \textbf{hybrid rule--LLM planner} that guarantees interpretable shot progression (e.g., pursuit $\rightarrow$ tracking shot $\rightarrow$ close-up $\rightarrow$ medium shot) while producing Chinese-first prompts optimized for models such as Doubao Seedream.
We further introduce an optional \textbf{Director Agent mode} with cross-shot memory, multi-candidate generation, heuristic and vision-based critics, edit-image continuity, and reflection logs.
The system is delivered as a deployable web application (React + FastAPI) with asynchronous tasks and server-sent event (SSE) progress streaming.
Experiments on chase, romance, and suspense prompts demonstrate coherent multi-shot plans and usable storyboard outputs.
\textbf{Code:} \url{https://github.com/ZiZuoLee/AI-Director-Agent}
\end{abstract}

\section{Introduction}
\label{sec:intro}

Storyboarding is a core pre-production step in filmmaking: directors translate screenplay beats into a sequence of framed shots before shooting or final rendering.
Recent text-to-image models can produce impressive single frames, but they lack (i)~\textit{shot decomposition}, (ii)~\textit{cinematographic reasoning}, and (iii)~\textit{inspectable intermediate representations}.
Course projects in computer graphics therefore benefit from systems that are not only generative, but also \textit{agentic} and \textit{explainable}.

We build \textbf{AI Director Agent} to address the following question:
\emph{How can a Chinese short story be automatically transformed into a multi-shot cinematic storyboard with transparent planning and optional director-level quality control?}

Our contributions are fourfold:
\begin{enumerate}[leftmargin=*, itemsep=2pt]
  \item A \textbf{hybrid rule--LLM planner} that combines interpretable cinematography rules with LLM-based Chinese prompt refinement.
  \item A \textbf{Director Agent loop} with memory, candidate search, heuristic and vision critics, and edit-reference continuity.
  \item A \textbf{deployable three-layer system} (Agent / Generation / System) exposed through FastAPI and a React web UI.
  \item An \textbf{explainable workflow} that surfaces parser tags, rule rationales, LLM JSON, and agent reflections to users.
\end{enumerate}

The minimum viable product (MVP) accepts one Chinese story, produces 3--5 shots, generates corresponding images, and merges them into a storyboard PNG.

\section{Related Work}
\label{sec:related}

\subsection{Text-to-Image and Creative Agents}
Diffusion models such as Latent Diffusion \citep{rombach2022high} enable high-quality image synthesis from natural-language prompts, but are typically used in a single-shot manner without film grammar or narrative progression.
Large language models \citep{brown2020language} support structured JSON planning; API aggregators such as OpenRouter \citep{openrouter2024} simplify access to multiple chat models.
Our planner builds on this stack for Chinese prompt refinement.

\subsection{Storyboard Pre-visualization}
Traditional storyboarding relies on manual framing decisions (shot scale, camera movement, emotional emphasis).
We encode a lightweight subset of these conventions as explicit rules rather than relying entirely on free-form LLM inference.

\subsection{Agentic Image Pipelines}
Recent systems combine tool use, memory, and critique loops for higher-quality media generation.
We adopt this pattern in Director Agent mode via cross-shot memory, candidate search, heuristic scoring, and multimodal vision analysis.
Implementation relies on FastAPI \citep{fastapi2024}, React \citep{react2025}, the Google GenAI SDK \citep{googlegenai2025}, Pillow \citep{pillow2025}, and ZenMux \citep{zenmux2025}.

\section{Method}
\label{sec:method}

\subsection{System Architecture}
The system follows a three-layer architecture aligned with our technical design document (TRD).
Figure~\ref{fig:pipeline} summarizes the end-to-end data flow.

\begin{figure}[t]
\centering
\fcolorbox{linegray}{blockblue}{%
  \begin{minipage}{0.9\linewidth}
    \centering
    \vspace{0.55em}
    {\bfseries Chinese Story Input}\\[0.35em]
    $\downarrow$\\[0.2em]
    {\bfseries Agent Layer}\\
    {\small \textit{parser $\rightarrow$ rules $\rightarrow$ llm $\rightarrow$ planner}}\\[0.35em]
    $\downarrow$\\[0.2em]
    {\bfseries Generation Layer}\\
    {\small \textit{prompt\_gen $\rightarrow$ generate / director\_agent}}\\[0.35em]
    $\downarrow$\\[0.2em]
    {\bfseries System Layer}\\
    {\small \textit{FastAPI orchestration + React UI + SSE progress}}\\[0.35em]
    $\downarrow$\\[0.2em]
    {\bfseries Storyboard PNG Output}\\[0.55em]
  \end{minipage}}
\caption{End-to-end pipeline of AI Director Agent.}
\label{fig:pipeline}
\end{figure}

\subsection{Agent Layer: Hybrid Shot Planning}

\textbf{Semantic parsing.}
Given input text $T$, \texttt{parser.py} segments sentences and extracts actions, emotions, locations, characters, and themes using bilingual keyword rules.
This step is intentionally lightweight and deterministic for reproducibility.

\textbf{Rule engine.}
\texttt{rules.py} maps semantic tags to cinematography templates.
For example, pursuit-related actions trigger a progression: \textit{tracking shot} $\rightarrow$ \textit{close-up} $\rightarrow$ \textit{medium shot}.
Camera movement (e.g., dolly, slow push-in, static) is derived from shot type and action cues.

\textbf{LLM prompt design.}
For each shot $i \in \{1,\dots,N\}$, the planner builds a Chinese seed prompt and calls an OpenRouter-compatible chat model \citep{openrouter2024,deepseek2025}.
The model returns JSON fields: \texttt{shot\_type}, \texttt{camera\_movement}, \texttt{prompt}, \texttt{description}, and \texttt{reason}.
We enforce Chinese-only diffusion prompts (50--120 characters), concrete character appearance, and per-shot narrative beats.
Shot~1's appearance description is reused as a \textit{character anchor} for later shots.
The planner output is $\mathcal{P} = \{\text{input\_text}, \text{parsed}, \text{shots}[1..N]\}$.

\subsection{Generation Layer}

\textbf{Prompt specification.}
\texttt{prompt\_gen.py} converts each planned shot into a backend-neutral \texttt{PromptSpec} with Chinese-first positive prompts, per-shot seeds, and 16:9 layout defaults.

\textbf{Simple mode.}
\texttt{generate\_images()} invokes the ZenMux backend with Doubao Seedream \citep{seedream2025} to produce one image per shot.

\textbf{Director Agent mode.}
\texttt{DirectorAgent} executes a per-shot quality-control loop (Algorithm~\ref{alg:director}).

\begin{algorithm}[t]
\caption{Director Agent per-shot loop}
\label{alg:director}
\begin{algorithmic}[1]
  \STATE Update cross-shot memory for characters and scenes
  \STATE Select strategy: \texttt{base}, \texttt{continuity\_anchor}, \texttt{scene\_anchor}, or \texttt{edit\_reference}
  \STATE Generate $K$ candidate images with seed offsets
  \STATE Score candidates with heuristic critic; re-rank top items with vision LLM
  \IF{best score $< \tau$ and attempts $< M$}
    \STATE Retry generation with adjusted strategy/prompt
  \ENDIF
  \STATE Log reflection (strategy, issue, action, score, rationale)
  \STATE Optionally edit from previous best frame for continuity
\end{algorithmic}
\end{algorithm}

\subsection{System Layer and Web UI}
\texttt{main.py} orchestrates planning, generation, and storyboard merging with progress callbacks.
\texttt{api.py} exposes REST endpoints and SSE events for asynchronous tasks.
\texttt{image\_merge.py} resizes shots, concatenates them horizontally, and exports \texttt{storyboard.png}.
The React UI provides script input, mode controls, and three inspection tabs: plan, images/storyboard, and agent logs.

\subsection{Design Highlights}
\label{sec:innovations}
\begin{itemize}[leftmargin=*, itemsep=2pt]
  \item \textbf{Hybrid rule--LLM planning} ensures film-language structure while LLMs refine visual prompts.
  \item \textbf{Chinese-first prompt chain} improves semantic alignment for Seedream-class backends.
  \item \textbf{Explainable artifacts} make parser tags, rule reasons, LLM JSON, and reflections user-visible.
  \item \textbf{Production-oriented delivery} replaces CLI-only scripts with async API, SSE, and a full web demo.
\end{itemize}

\section{Experiments and Results}
\label{sec:experiments}

\subsection{Experimental Setup}
We evaluate three representative Chinese prompts with a fixed generation configuration (Table~\ref{tab:setup}).

\begin{table}[t]
\centering
\small
\begin{tabular}{@{}ll@{}}
\toprule
\textbf{Setting} & \textbf{Value} \\
\midrule
Test scenarios & Pursuit, Romance, Suspense \\
Shots per story & $N = 3$ \\
Resolution / aspect & $1024 \times 576$, 16:9 \\
Planning model & OpenRouter DeepSeek \\
Image model & ZenMux Doubao Seedream \\
Generation modes & Simple and Director Agent \\
\bottomrule
\end{tabular}
\caption{Experimental configuration.}
\label{tab:setup}
\end{table}

\textbf{Pursuit.} A protagonist is chased through a dark alley.
\textbf{Romance.} A campus rooftop confession at dusk.
\textbf{Suspense.} A detective discovers footprints in an abandoned factory.

\subsection{Qualitative Results}
The hybrid planner produces distinct per-shot intents (e.g., tracking $\rightarrow$ emotional close-up $\rightarrow$ confrontation medium shot) and Chinese prompts with concrete appearance attributes.
Simple mode yields usable cinematic frames and a merged storyboard preview in the UI.
Director mode improves cross-shot consistency when edit-reference is triggered, at higher latency.
Users can inspect parser tags, per-shot prompts, generated images, and (in agentic mode) candidate rankings and reflection logs.

\subsection{Observations}
Table~\ref{tab:compare} summarizes mode-level trade-offs observed in our runs.

\begin{table}[t]
\centering
\small
\begin{tabular}{@{}lcc@{}}
\toprule
\textbf{Aspect} & \textbf{Simple Mode} & \textbf{Director Mode} \\
\midrule
Latency & Low & 3--5$\times$ slower \\
Cross-shot consistency & Prompt anchoring & Memory + edit-reference \\
Explainability & Plan JSON only & Scores + reflection logs \\
Best use case & Fast demo / MVP & Higher-quality continuity \\
\bottomrule
\end{tabular}
\caption{Comparison between generation modes.}
\label{tab:compare}
\end{table}

Rule-only prompts are less visually specific than LLM-refined prompts.
Chinese-first prompt assembly improves semantic alignment for Seedream compared with English suffixes.

\subsection{Limitations}
\begin{itemize}[leftmargin=*, itemsep=2pt]
  \item Parser is keyword-based rather than LLM-semantic parsing.
  \item Character consistency relies on prompt anchoring and optional edit-reference.
  \item Task state is in-memory and not persisted across server restarts.
  \item Image backend is API-based (ZenMux) rather than local diffusion pipelines.
\end{itemize}

\section{Conclusion}
\label{sec:conclusion}

We presented AI Director Agent, a deployable Chinese storyboard generation system that combines interpretable cinematography rules, LLM prompt design, image synthesis, optional director-level critique loops, and storyboard composition in a unified web workflow.
The project demonstrates how agentic decomposition and explainable intermediate representations can elevate text-to-image from single-frame generation to structured cinematic pre-visualization.
Future work includes LLM-based screenplay parsing, plan-then-edit user workflows, persistent task queues, additional image backends, and video/animatic export.

\section*{Team Contributions}
\label{sec:team}

Division of labor follows the repository TRD (Agent / Generation / System layers) and is corroborated by Git commit history.

\begin{table}[t]
\centering
\footnotesize
\setlength{\tabcolsep}{4pt}
\renewcommand{\arraystretch}{1.15}
\begin{tabularx}{\linewidth}{@{}l l l X@{}}
\toprule
\textbf{Name} & \textbf{ID} & \textbf{Role} & \textbf{Contributions} \\
\midrule
Lee Zi Zuo (Lv Zizhuo) & 23300246003 & Agent &
PRD/TRD; parser; rule engine; planner; OpenRouter LLM workflow. \\
Thai Tian Qi (Dai Tianqi) & 23300246006 & Generation &
ZenMux backend; prompt spec; Director Agent; memory and critic modules. \\
Carole Hew (Qiu Jiakun) & 23300246008 & System &
FastAPI + SSE; React UI; pipeline orchestration; storyboard merge. \\
\bottomrule
\end{tabularx}
\caption{Team roster and division of labor.}
\label{tab:team}
\end{table}

\bibliography{references}
\bibliographystyle{iclr2025_conference}

\end{document}
